% =====================================================================================
% Chapter 4 · Segment effects — pooling, shrinkage, and the price of the average
%
% Built from notebooks/02_segment_effects.ipynb. NOT ONE NUMBER IN THIS FILE IS TYPED BY
% HAND: every figure in the prose is a macro from build/macros.tex, generated by cmp.report
% from the executed notebook. The only literals below are (i) the constants of the
% data-generating model in §4.2, which are *definitions* rather than results, and (ii) the
% decision-rule thresholds (0.8 / 0.5) and the 0.05 test level, which are *conventions*
% stated in the text, not measurements.
% =====================================================================================
\chapter{Segment effects: pooling and shrinkage}
\label{ch:seg}

\begin{quote}\small
\emph{An email campaign is sent to a randomized half of a \nbTwoN-customer base at a contact
cost of \eur{\nbTwoCost} a head. Pooled over everyone it lifts spend by \eur{\nbTwoCpEst}, which
clears the cost comfortably, and the pooled recommendation is therefore to mail the entire base.
That recommendation is wrong for every customer in it. The effect is genuinely different across
tiers --- \eur{\nbTwoTrueLow} for low-value customers and \eur{\nbTwoTrueHigh} for high-value ones
by construction --- so the pooled estimate misses the low tier by \eur{\nbTwoCpMissLow} and mails
\pc{\nbTwoShareLow} of the base at a loss, while reporting the tightest interval on the page. The
repair is not a bigger model but an older idea: shrink each segment's own estimate toward the grand
mean by a weight the data choose. Pointed at \nbTwoNNoise{} segments manufactured from coin flips,
that weight collapses to $\lambda = \nbTwoNoiseLam$ and every sham ``responsive segment'' is
annihilated; pointed at the two real tiers it returns $\lambda = \nbTwoRealLam$ and leaves them
alone. And the Bayesian machinery this book is nominally selling buys, on this problem,
\emph{nothing} on location --- the classical and Bayesian estimates agree to \eur{\nbTwoMaxPairGap}
--- and nothing on width. What it buys is \cref{sec:seg:euro}'s send rule, and
\cref{sec:seg:compare} says so plainly.}
\end{quote}

% =====================================================================================
\section{The decision}
\label{sec:seg:decision}

A retailer emails an offer to a randomized half of its customer base. The campaign works: averaged
over everybody, emailed customers spend more, and the average lift beats the \eur{\nbTwoCost} it
costs to contact a customer. The average is also the least interesting number in the file.

The question the marketing team actually has is \emph{who} the email works on. Suppose it moves
high-value, highly engaged customers a great deal and barely touches everyone else. Then the
average is a blend of a large effect and a small one, and acting on the blend does two expensive
things at once: it mails customers for whom the offer does not pay, and --- when the cost rises, or
the margin thins --- it eventually stops mailing the customers for whom it pays handsomely, because
the \emph{average} has stopped clearing the bar. A single averaged decision has no way to keep the
profitable segment.

The mirror-image error is at least as expensive and much easier to commit. Slice the base finely
enough and some slice will look responsive purely by chance; write the send rule around that slice
and the firm has fitted its marketing to noise. \Cref{sec:seg:wrong} manufactures such a discovery
on this chapter's own data, from a coin flip.

So the decision is not one number but a bracket, and the chapter is organised around it:

\begin{description}[leftmargin=1.4em, style=nextline, itemsep=.3em]
  \item[No pooling.] Estimate each segment from its own rows alone. Nothing is borrowed and nothing
    is assumed --- which is exactly why a thin segment comes back wild, with an interval wide enough
    to accommodate any policy at all.
  \item[Complete pooling.] One number for everybody. It is the most precise estimate available, and
    if the segments genuinely differ it is wrong for every one of them at once. Precision around
    the wrong number is the more dangerous failure, because it looks like confidence.
  \item[Partial pooling.] The interpolation between the two --- and, crucially, an
    \emph{adaptive} one. It shrinks a thin, noisy segment hard toward the grand mean and barely
    touches a fat, well-measured one, with the amount of shrinkage chosen by the data rather than
    by the analyst.
\end{description}

Everything below lives in the gap between the first two, and the third is what fills it.

% =====================================================================================
\section{The data-generating model}
\label{sec:seg:dgm}

A method that claims to find heterogeneity must first be shown not to invent it, and that
demonstration requires a world whose heterogeneity is known. The simulator draws
$i = 1, \dots, \nbTwoN$ customers. Each carries a value tier $v_i \in \{0, 1\}$ (\emph{high-value}
$= 1$; \nbTwoNHigh{} of the base, against \nbTwoNLow{} low-value), a continuous engagement score
$g_i \sim U(0,1)$, and a background \emph{trend} covariate $t_i$ that also drives spend. The email
$e_i \in \{0,1\}$ is assigned by a fair coin, independently of everything else.

The planted per-customer treatment effect is linear in the two moderators and in nothing else:
%
\begin{equation}
  \tau_i \;=\; 3 \;+\; 9\,v_i \;+\; 8\,g_i ,
  \label{eq:seg:tau}
\end{equation}
%
and spend is that effect, switched on by the email, plus a baseline that depends on the same
customer characteristics:
%
\begin{equation}
  Y_i \;=\; \beta_0 \;+\; \tau_i\, e_i \;+\; \beta_2 v_i \;+\; \beta_5 g_i \;+\; \beta_6 t_i
            \;+\; \varepsilon_i,
  \qquad \varepsilon_i \sim \mathcal N\!\left(0,\ 6^{2}\right).
  \label{eq:seg:dgp}
\end{equation}
%
Three features of \cref{eq:seg:tau} carry the chapter. First, the heterogeneity is \emph{real and
large}: the effect runs from \eur{\nbTwoTauMin} for a disengaged low-value customer to
\eur{\nbTwoTauMax} for a fully engaged high-value one. Second, it is \emph{layered} --- a discrete
tier gap of €9 and a continuous engagement slope of €8 --- so a model that bins customers into two
groups and stops there is leaving structure on the table. Third, at the average engagement
$g = 0.5$ the two tier effects work out to \eur{\nbTwoTrueLow} (low) and \eur{\nbTwoTrueHigh}
(high), and \emph{those} are the numbers every estimator in this chapter is graded against. The
population average effect --- what a pooled analysis is honestly estimating --- is
\eur{\nbTwoTrueAte}.

The contact cost is \eur{\nbTwoCost}, and it is placed deliberately: it sits \emph{between} the two
planted segment effects. A campaign that cost €2 would be worth running on everybody and a campaign
that cost €20 on nobody, and in either world segmentation would be an academic exercise. At
\eur{\nbTwoCost} the low tier is a marginal call and the high tier is not, so the pooling choice
changes what the firm does with \pc{\nbTwoShareLow} of its base. That is the whole point of the
number.

% =====================================================================================
\section{Identification}
\label{sec:seg:ident}

\subsection*{The estimand}

Write $Y_i(1)$ and $Y_i(0)$ for the spend of customer $i$ with and without the email
\citep{rubin1974, imbens2015}. Only one of the two is ever observed. When the effect depends on who
receives the treatment, the effect is said to be \emph{moderated}, and the characteristic it depends
on is a \emph{moderator}. The object of interest is then not the average effect but the
\emph{conditional} average treatment effect --- the CATE --- the average effect among customers who
look like $x$:
%
\begin{equation}
  \tau(x) \;=\; \E\big[\,Y_i(1) - Y_i(0) \;\big|\; X_i = x\,\big],
  \qquad
  \tau \;=\; \E\big[\tau(X)\big] .
  \label{eq:seg:cate}
\end{equation}
%
The ATE of \cref{eq:seg:cate} is the CATE averaged over the base; a segment effect $\tau_s$ is the
CATE averaged over segment $s$. All three are the same object read at different resolutions, which
is why the arms of this chapter can be compared at all: they estimate one estimand, not several.
Throughout, the two segment effects are read \emph{at mean engagement}, so that the quantity on the
page is the planted \eur{\nbTwoTrueLow} and \eur{\nbTwoTrueHigh} and nothing else.

\subsection*{Where heterogeneity lives}

Suppose spend is generated by a structural equation linear in the parameters, with the email
interacted with the two moderators:
%
\begin{equation}
  \E[\,Y \mid e, v, g, t\,] \;=\; \beta_0 + \beta_1 e + \beta_2 v + \beta_3\,(e \cdot v)
                                  + \beta_4\,(e \cdot g) + \beta_5 g + \beta_6 t .
  \label{eq:seg:structural}
\end{equation}
%
Intervene on the email --- set $e=1$, then $e=0$, in \citeauthor{pearl2009}'s $do(\cdot)$ sense
\citep{pearl2009} --- and difference. Every term not multiplied by $e$ cancels, and what survives is
%
\begin{equation}
  \tau(x) \;=\; \E[Y \mid do(e{=}1), x] - \E[Y \mid do(e{=}0), x]
          \;=\; \beta_1 \;+\; \beta_3\, v \;+\; \beta_4\, g .
  \label{eq:seg:tauhat}
\end{equation}
%
\Cref{eq:seg:tauhat} is the algebraic spine of the chapter. \textbf{All of the heterogeneity is
carried by the interaction coefficients.} If $\beta_3 = \beta_4 = 0$ the effect is the same for
everyone and \cref{eq:seg:cate} collapses to the ATE. So the question ``is the effect
heterogeneous?'' is not a vague one about segments; it is the sharp question ``is $\beta_3$
(or $\beta_4$) non-zero?'', and it is answered by an interval on a coefficient.

\subsection*{The assumptions}

\begin{assumption}[Unconfoundedness]\label{as:seg:unconf}
$\big(Y(1), Y(0)\big) \perp\!\!\!\perp e \mid X$. Here it holds \emph{by design}: the email was
assigned by a coin flip, so nothing points into $e$ and the admissible adjustment set the software
reads off the graph is the empty set. This is the assumption that does all the causal work in the
chapter, and it is the one that is free here and expensive on real data --- where an email is sent
to whoever the CRM thought would respond, and the segment effects are then contaminated by
\emph{who got emailed}. \Cref{ch:seg}'s machinery is unchanged in that world; only this assumption
is, and it is not a small change.
\end{assumption}

\begin{assumption}[Positivity within every segment]\label{as:seg:pos}
$0 < \Prob(e = 1 \mid X = x) < 1$ for every $x$ on which a CATE is read. A segment in which nobody
was emailed --- or everybody was --- yields no comparison, only an extrapolation. It is
\emph{testable and tested}: the emailed share is near one half in both value tiers and in every
engagement tercile, so each segment's CATE is a comparison of two arms that are both in force. On
real CRM data this is the check that quietly voids a segment, and it must be run before, not after,
the model.
\end{assumption}

\begin{assumption}[No interference]\label{as:seg:sutva}
Emailing customer $i$ does not change customer $j$'s spend --- SUTVA \citep{imbens2015}. Household
members share a basket, and a forwarded offer code is interference in its purest form. \textbf{No
test in this chapter can see it}, and none is offered: it is discharged by design (deduplicate
households, one code per customer) or it is not discharged at all. An assumption with no
accompanying test is a wish; the honest move is to say which of the two it is.
\end{assumption}

\begin{assumption}[The moderation form]\label{as:seg:form}
The effect is linear in the moderators, as in \cref{eq:seg:tauhat}. This assumption is
\emph{additional} to identification --- \cref{as:seg:unconf} identifies $\tau(x)$ without it --- and
it is the one that the simulator quietly hands the analyst and reality does not. It is tested three
ways in \cref{sec:seg:valid,sec:seg:wrong}: by a posterior predictive check, by a deliberately
misspecified refit, and by a model-free cross-check that assumes no functional form at all. The
graph-falsification test, notably, \emph{cannot} test it --- \cref{sec:seg:wrong} shows why.
\end{assumption}

\begin{assumption}[Exchangeable segments]\label{as:seg:exch}
For the partial-pooling estimator of \cref{sec:seg:classical} only: the segment effects are draws
from a common distribution, $\tau_s \sim \mathcal N(\bar\tau, \sigma_\tau^2)$, with a spread
$\sigma_\tau$ that the data estimate rather than the analyst assert. This is the assumption that
licenses one segment to borrow strength from another, and it is the only assumption in this chapter
that \emph{buys} something rather than merely permitting it. It is also the one most often smuggled
in: pooling across segments that are not exchangeable (a test market and a control market; a
product line being discontinued) imports bias under the banner of borrowing strength. It is tested
directly in \cref{sec:seg:classical}, where the estimated $\sigma_\tau$ is shown to go to zero on
segments that are noise and to \eur{\nbTwoRealSigma} on segments that are real.
\end{assumption}

% =====================================================================================
\section{The classical baseline}
\label{sec:seg:classical}

No likelihood, no prior, no sampler. What does a competent analyst produce here without any of them?

\subsection*{The two ends of the bracket}

The classical toolkit's first two moves are the two extremes of \cref{sec:seg:decision}. Writing
$\tau_s$ for segment $s$'s effect, $n_s$ for its row count and $\hat\tau_s$ for an estimate:
%
\begin{equation}
  \underbrace{\hat\tau^{\,\mathrm{np}}_s \;=\; \text{OLS on the } n_s \text{ rows of segment } s
              \text{ alone}}_{\textbf{no pooling: unbiased, variance} \;\propto\; 1/n_s}
  \qquad\qquad
  \underbrace{\hat\tau^{\,\mathrm{cp}}_s \;\equiv\; \hat\tau \quad \forall s}
             _{\textbf{complete pooling: precise, biased}}
  \label{eq:seg:bracket}
\end{equation}
%
Both arms target the \emph{same} per-segment estimand under the \emph{same} identification. What
differs is only how much information each is willing to borrow across segments: none, or all of it.

\Cref{tab:seg:classical} runs both. No pooling puts each segment within a euro of its planted truth,
with intervals \eur{\nbTwoNpLowWidth} and \eur{\nbTwoNpHighWidth} wide --- with \nbTwoNLow{}
low-value and \nbTwoNHigh{} high-value customers, splitting the sample costs almost nothing, because
a few hundred rows per arm is already plenty. It is worth pausing on the coverage column, which does
not come out the way a tidy textbook would want: the high-value interval covers its truth and the
low-value interval \nbTwoNpLowCovers{} it, closing a hair above \eur{\nbTwoTrueLow}. A nominal
90\,\% interval is \emph{supposed} to miss one time in ten, and here is one of the misses on the
first table in the chapter. The honest response is to point at it rather than re-roll the seed until
the column reads all-yes --- and to note that this sample's \eur{\nbTwoNpErrLow} of upward error on
the low tier will reappear in \emph{every} arm below that estimates that segment separately,
classical and Bayesian alike, because it is a property of the sample, and no estimator argues with
the sample it was given.

\input{build/tables/nb02_classical}

Complete pooling is the interesting failure. It returns \eur{\nbTwoCpEst} with the tightest interval
on the page --- \pc{\nbTwoCpNarrowerPct} narrower than the no-pooling average --- and that interval
excludes \emph{both} truths. It is not a broken estimator. It is an excellent estimator of a
quantity nobody asked for: it lands squarely on the planted ATE of \eur{\nbTwoTrueAte}, and it
answers the wrong question with impressive confidence, which is the most expensive way to be wrong.

The expense is immediate and countable. At a contact cost of \eur{\nbTwoCost} the pooled estimate of
\eur{\nbTwoCpEst} says \textsc{mail everybody}. The low tier's own effect --- planted at
\eur{\nbTwoTrueLow}, estimated at \eur{\nbTwoNpLow} --- does not clear the cost. Pooling therefore
does not merely blur the picture; it mails the \pc{\nbTwoShareLow} of the base that loses money, and
it does so while reporting the narrowest confidence interval in the analysis.

\subsection*{The covariance choice, and the interval that is not a probability}

The standard errors above are heteroskedasticity-robust (HC1; \citealp{white1980, mackinnon1985}).
These are cross-sections of independent customers --- no panel, no repeated measurement, no time
series --- so there is no cluster to cluster on, and the \citet{liang1986} machinery that a panel
would demand has nothing here to grip. The honest default is robustness to
\emph{heteroskedasticity}, not to correlation.

Rather than assert that the choice is harmless, the notebook refits both and prints the two side by
side. On the low tier the robust standard error is \eur{\nbTwoSeRobustLow} against
\eur{\nbTwoSeIidLow} for the iid formula; on the high tier, \eur{\nbTwoSeRobustHigh} against
\eur{\nbTwoSeIidHigh}. They agree to within \pc{\nbTwoSeGapPct}. That is worth stating plainly,
because it is a small result with a large moral: on a well-behaved simulator that draws
homoskedastic noise for everybody, \textbf{the robust sandwich costs nothing and buys nothing}. It
is insurance, not magic. On real spend data the premium would not be zero --- high-value customers
spend more \emph{and vary more}, and an iid standard error would then average the two tiers' noise
levels into one, reporting a too-narrow interval for the noisy segment and a too-wide one for the
quiet one. Take the insurance anyway. It is free when you do not need it, which is the definition of
good insurance.

One thing the table cannot do, at any level of effort, is attach a probability to a segment's
effect. The low tier's 90\,\% interval, $[\eur{\nbTwoNpLowLo}, \eur{\nbTwoNpLowHi}]$, straddles the
\eur{\nbTwoCost} contact cost. The classical apparatus can therefore say \emph{``we cannot reject
that low-value customers break even''} --- and it can say nothing else. That sentence is exactly as
compatible with a segment 20\,\% likely to clear its cost as with one 60\,\% likely, and it will not
say which, although those two worlds call for opposite campaigns. \Cref{sec:seg:euro} needs
$\Prob(\tau_{\mathrm{low}} > \eur{\nbTwoCost})$, a number between zero and one about \emph{this
segment's effect}. It does not exist in the frequentist vocabulary, because there the effect is not
a random variable. That, and not a better point estimate, is what \cref{sec:seg:bayes} is for.

\subsection*{Partial pooling, without a sampler: the shrinkage factor}

The middle of the bracket is not a Bayesian idea and it does not require one. It is
\citeauthor{stein1956}'s \citeyearpar{stein1956}, and its practical form is the empirical-Bayes or
James--Stein estimator \citep{james1961, efron1975}. Given per-segment estimates $\hat\tau_s$ with
standard errors $\hat{\mathrm{se}}_s$, pull each toward the grand mean $\bar\tau$ by a weight the
data choose:
%
\begin{equation}
  \tilde\tau_s \;=\; \lambda_s\,\hat\tau_s \;+\; (1 - \lambda_s)\,\bar\tau,
  \qquad
  \lambda_s \;=\; \frac{\hat\sigma^2_\tau}{\hat\sigma^2_\tau + \hat{\mathrm{se}}_s^2},
  \qquad
  \hat\sigma^2_\tau \;=\; \max\!\Big(0,\ \widehat{\Var}(\hat\tau)
                                        - \overline{\hat{\mathrm{se}}^2}\Big).
  \label{eq:seg:shrink}
\end{equation}
%
Every term earns its place. $\hat\sigma^2_\tau$ is the \emph{between-segment} spread that survives
after the sampling noise is subtracted off: the observed scatter of the estimates,
$\widehat{\Var}(\hat\tau)$, is inflated by the fact that each $\hat\tau_s$ is measured with error,
so that error is netted out --- and the $\max(0, \cdot)$ records the possibility that once it is,
nothing is left.

The reason \cref{eq:seg:shrink} is a principled compromise rather than a fudge is best seen by
rewriting it. Divide numerator and denominator by $\hat\sigma^2_\tau\,\hat{\mathrm{se}}_s^2$:
%
\begin{equation}
  \tilde\tau_s \;=\;
    \frac{\dfrac{1}{\hat{\mathrm{se}}_s^{2}}\ \hat\tau_s \;+\;
          \dfrac{1}{\hat\sigma^{2}_\tau}\ \bar\tau}
         {\dfrac{1}{\hat{\mathrm{se}}_s^{2}} \;+\; \dfrac{1}{\hat\sigma^{2}_\tau}} .
  \label{eq:seg:precision}
\end{equation}
%
\Cref{eq:seg:precision} is a \textbf{precision-weighted average of two estimates of the same
quantity}: the segment's own noisy estimate, weighted by its own precision, and the grand mean,
weighted by the precision with which the population of segments is concentrated. It is the same
arithmetic that combines two measurements of a physical constant, and it is optimal for the same
reason. Nothing about it is Bayesian, although the Bayesian version --- which replaces the plug-in
$\hat\sigma_\tau$ with a posterior over $\sigma_\tau$, and thereby propagates the uncertainty in the
amount of borrowing --- is a short walk away.

Read $\lambda_s$ as \emph{``how much of segment $s$'s own estimate do we keep?''} and the two limits
of \cref{eq:seg:bracket} fall out of one formula:
%
\begin{equation}
  \hat\sigma^2_\tau \gg \hat{\mathrm{se}}_s^2 \;\Longrightarrow\; \lambda_s \to 1
     \quad (\text{no pooling}),
  \qquad\qquad
  \hat\sigma^2_\tau \to 0 \;\Longrightarrow\; \lambda_s \to 0
     \quad (\text{complete pooling}).
  \label{eq:seg:limits}
\end{equation}
%
\textbf{The two classical arms are the two limits of a single estimator, and $\lambda$ decides where
between them you sit.} It is a signal-to-noise ratio: the ratio of how different the segments truly
are to how badly each one is measured. And because $\lambda_s$ carries the subscript, the shrinkage
is \emph{adaptive across segments as well as across problems} --- a thin, noisy segment with a large
$\hat{\mathrm{se}}_s$ is pulled hard toward the grand mean while a fat one sitting beside it is
barely touched. That is the property no fixed rule has.

\subsection*{$\lambda$ is learned: the demonstration}

The claim that $\lambda$ is chosen by the data rather than by the analyst is checkable, and the
notebook checks it by pointing the identical estimator at two kinds of segment
(\cref{fig:seg:shrink}).

First, at \nbTwoNNoise{} \textbf{pure-noise segments}: the same customers, the same randomized
email, the same spend --- but labels assigned by a coin flip, so that the true effect is
\emph{identical} in every one of them and $\sigma_\tau$ is exactly zero by construction. No-pooling
hands back \nbTwoNNoise{} different numbers spread over
\eur{\nbTwoNoiseRawLo}--\eur{\nbTwoNoiseRawHi} --- it can do nothing else; it will always hand back
as many numbers as there are labels, and some of them will look impressive. \Cref{eq:seg:shrink}
recovers $\hat\sigma_\tau = \eur{\nbTwoNoiseSigma}$, sets $\lambda = \nbTwoNoiseLam$, and collapses
the entire range onto the grand mean \eur{\nbTwoNoiseGrand}. \textbf{Partial pooling discovers that
there is no heterogeneity to find} --- which is a conclusion no-pooling is structurally incapable of
reaching, and it is the single most useful thing this estimator does.

Second, at the \textbf{two real value tiers}, whose effects genuinely differ by €9. Here
\cref{eq:seg:shrink} recovers $\hat\sigma_\tau = \eur{\nbTwoRealSigma}$, sets
$\lambda = \nbTwoRealLam$, and leaves the estimates essentially untouched. \Cref{as:seg:exch} is
therefore not a wish but a measured quantity, and it is measured in the right direction both times.
Shrinkage is not a haircut applied to everything. It is a signal-to-noise ratio, and on fat,
truly-different segments it declines to shrink.

\input{build/floats/nb02_shrink}

\subsection*{When does the pooling choice bite?}

Two fat segments is the friendly case, and honesty requires the chapter to say what its own headline
method is worth there: not much. The notebook therefore cuts the same base into \nbTwoSweepCells{}
real cells --- value tier $\times$ engagement sextile --- thins them, and runs all three estimators
against each cell's own true effect over \nbTwoSweepReps{} fresh populations
(\cref{tab:seg:sweep}).

Three things happen, and they are the bias--variance bracket made numerical.

\begin{enumerate}[leftmargin=*]
  \item \textbf{No pooling degrades as the cells thin,} exactly as $1/n_s$ predicts: its RMSE grows
        from \eur{\nbTwoSweepFatNp} at \nbTwoSweepFatRows{} rows per cell to \eur{\nbTwoSweepThinNp}
        at \nbTwoSweepThinRows, and its interval width grows
        \nbTwoSweepWidthGrowth$\times$, from \eur{\nbTwoSweepFatWidth} to
        \eur{\nbTwoSweepThinWidth}. An interval that wide accommodates any policy at all, which is
        another way of saying it supports none.
  \item \textbf{Complete pooling's error is flat and large}
        (\eur{\nbTwoSweepFatCp} against \eur{\nbTwoSweepThinCp}). It does not improve as the cells
        thin, and it would not improve if the cells were fat: that is \emph{bias}, and no amount of
        data removes it. This is the single most important row of the table, because it is the row
        that explains why the pooled estimate above could be both the most precise and the most
        wrong.
  \item \textbf{Shrinkage tracks the better of the two at every size,} and the learned $\lambda$
        falls from \nbTwoSweepFatLam{} to \nbTwoSweepThinLam{} as the cells thin --- borrowing more
        precisely when there is less to go on.
\end{enumerate}

\input{build/tables/nb02_sweep}

But read the $\lambda$ column before celebrating. It stays high, so the gain over plain no pooling
is real but modest --- \eur{\nbTwoSweepThinSh} against \eur{\nbTwoSweepThinNp} at the thinnest cell
size. The reason is visible in \cref{eq:seg:shrink}: the planted heterogeneity here is enormous, so
$\hat\sigma^2_\tau$ swamps $\hat{\mathrm{se}}_s^2$ and there is little for shrinkage to do.
\textbf{Partial pooling is not a free precision win, and any book that sells it as one is selling
something.} It is an adaptive guard whose value scales with how noisy the segments are relative to
how different they truly are --- which is precisely why the same estimator, pointed at the coin-flip
labels, drove $\lambda$ to \nbTwoNoiseLam{} and annihilated every sham discovery. It earns its keep
where the segments are thin and the effects are close. It earns very little here, and the chapter
says so.

Notice, finally, what has \emph{not} happened. No likelihood was written down, no prior was chosen,
no sampler ran. Randomization did the causal work, two OLS fits did the estimation, and a formula
from 1956 did the pooling. The Bayesian section must now justify itself against \emph{this}.

% =====================================================================================
\section{The Bayesian model}
\label{sec:seg:bayes}

The Bayesian arm fits \cref{eq:seg:structural} whole, in one model, rather than as two separate
regressions. Writing $i$ for the customer:
%
\begin{align}
  Y_i \mid \beta, \sigma &\;\sim\; \mathcal N\big(\mu_i,\ \sigma^{2}\big), \label{eq:seg:lik}\\
  \mu_i &\;=\; \beta_0 + \beta_1 e_i + \beta_2 v_i + \beta_3\, e_i v_i + \beta_4\, e_i g_i
               + \beta_5 g_i + \beta_6 t_i, \label{eq:seg:mu}\\
  \beta_k &\;\sim\; \mathcal N\!\left(0,\ 10^{2}\right), \qquad k = 0, \dots, 6,
               \label{eq:seg:prior_b}\\
  \sigma &\;\sim\; \text{HalfNormal}(1). \label{eq:seg:prior_s}
\end{align}
%
Inference is by the No-U-Turn sampler \citep{hoffman2014} in PyMC \citep{salvatier2016}:
\nbTwoNDraws{} post-warmup draws, $\hat R = \nbTwoRhat$, a minimum bulk effective sample size of
\nbTwoEss{} \citep{vehtari2021}, and \nbTwoDivergences{} divergent transitions. Nothing that follows
is a sampler artefact.

\subsection*{Are these priors weak?}

The question is malformed as usually asked, and the model above is a good place to say why. A prior
is weak or strong only \emph{relative to the precision of the likelihood}, never relative to the raw
scale of the data. \Cref{eq:seg:prior_s} looks alarming --- a HalfNormal(1) on a noise level that is
truly €6 --- and it is harmless: with \nbTwoN{} rows the likelihood is so sharp that the posterior
for $\sigma$ lands at \eur{\nbTwoSigmaPost}, where the data put it. The prior wanted something much
smaller and was simply overruled.

Reversing the experiment makes the same point in the direction that matters. Refit with
$\beta_k \sim \mathcal N(0, 0.1^2)$ --- a prior roughly an order of magnitude \emph{more} precise
than the likelihood, judged against the posterior standard deviations of the coefficients in the
default fit --- and the high-value CATE collapses from \eur{\nbTwoDefaultPriorHigh} to
\eur{\nbTwoTightPriorHigh}. The posterior is always prior $\times$ likelihood. The defaults were
harmless here \emph{only because $n$ made the likelihood sharp}, and the moment segments get thin
the per-segment likelihood weakens and the prior starts to bite. That is not a footnote; it is the
same $\lambda$ of \cref{eq:seg:shrink} wearing different clothes. \textbf{When segments are thin,
priors must be chosen, not defaulted} --- and the principled way to choose them is exactly the
hierarchical one: let the segments themselves estimate the prior's spread.

\subsection*{Three models, one estimand}

\Cref{fig:seg:forest} puts \cref{sec:seg:classical}'s bracket on a single axis, fitted Bayesianly:
full pooling (no interaction terms), no pooling (a separate fit inside each tier), and the
interaction model of \cref{eq:seg:structural}. Full pooling lands at \eur{\nbTwoBayesPool}, between
the two truths and therefore wrong for both segments at once --- too optimistic about low-value
customers and too pessimistic about high-value ones. No pooling is centred but re-estimates
$\sigma$, the engagement slope and the trend slope inside each half from its own rows.

\input{build/floats/nb02_forest}

The interaction model wins, and it is worth being precise about \emph{what} it wins, because it is
not what the brochure promises. At these segment sizes the intervals of (b) and (c) are nearly the
same width; sharing the nuisance parameters buys almost nothing. What (c) buys is \textbf{coherence}:
a single fit yields both segment effects, \emph{and} the posterior of their difference
($\beta_3$ --- the question ``is the gap real?''), \emph{and} the continuous engagement profile
($\beta_4$), \emph{and} --- the part with no classical closed form at all --- the break-even
engagement of \cref{sec:seg:euro}, a ratio of two estimated coefficients whose uncertainty falls
straight out of the joint posterior. Model (b) can only bolt the first of those together from two
independent posteriors, and model (a) denies the rest exist.

% =====================================================================================
\section{Diagnostics and validation}
\label{sec:seg:valid}

\subsection*{Can the model regenerate its own data?}

Before reading effects off a structural model, ask whether it can reproduce the outcome it was fit
to. Replicated spend vectors drawn from the fitted model are overlaid on the observed distribution
in \cref{fig:seg:ppc}: \pc{\nbTwoPpcCov} of observed spend falls inside its own replicates'
5--95\,\% band, against a nominal 90\,\%. No gross misfit. A model that cannot regenerate the
outcome's shape has no business issuing counterfactuals about it.

\input{build/floats/nb02_ppc}

\subsection*{Recovery, and the shape of the effect}

\Cref{fig:seg:recovery} grades the fit against the planted truth on three levels at once. The two
segment posteriors sit close to \eur{\nbTwoTrueLow} and \eur{\nbTwoTrueHigh} and barely overlap:
\eur{\nbTwoCateLow} (90\,\% HDI $[\eur{\nbTwoCateLowLo}, \eur{\nbTwoCateLowHi}]$) and
\eur{\nbTwoCateHigh} ($[\eur{\nbTwoCateHighLo}, \eur{\nbTwoCateHighHi}]$), against a pooled ATE of
\eur{\nbTwoAte}. The continuous profile recovers the planted engagement slope --- $\beta_4$ has
posterior mean \eur{\nbTwoBeg} against a planted \eur{\nbTwoBegTrue}, with a 90\,\% credible interval
of $[\eur{\nbTwoBegLo}, \eur{\nbTwoBegHi}]$, which \nbTwoBegCovers{} the truth on this draw. And the
per-customer fit hugs the 45-degree line, at an RMSE of \eur{\nbTwoRmseInd} on effects spanning
\eur{\nbTwoTauMin}--\eur{\nbTwoTauMax}: recovery at the individual level, not merely on segment
averages, which is the harder test and the one that licenses \cref{sec:seg:euro}'s per-customer send
rule.

\input{build/floats/nb02_recovery}

The gap itself is the heterogeneity verdict. The interaction coefficient $\beta_3$ has posterior
mean \eur{\nbTwoBev} against a planted \eur{\nbTwoBevTrue}, with $\Prob(\beta_3 > 0) =
\nbTwoPGap$. The segments truly differ, and the model says so.

\subsection*{One sample is one sample}

A single draw of the world can land a little off --- the low-value posterior above runs
\eur{\nbTwoSampleErrLow} high --- so the estimator is re-run on \nbTwoNSeeds{} fresh samples. The bias washes out to
\nbTwoSeedBiasLow{} euros on the low tier and \nbTwoSeedBiasHigh{} on the high tier, with
seed-to-seed standard deviations of \eur{\nbTwoSeedSdLow} and \eur{\nbTwoSeedSdHigh}, and the nominal
90\,\% intervals cover the truth on \nbTwoSeedCovLow{} and \nbTwoSeedCovHigh{} of \nbTwoNSeeds{}
samples respectively. The one-sample miss in \cref{tab:seg:classical} was the one-in-ten miss a
90\,\% interval is entitled to, and not a symptom.

The same sweep tracks a quantity \cref{sec:seg:euro} will lean on and which deserves separate
scrutiny --- the break-even engagement $g^\star$, a \emph{ratio} of estimated coefficients. Ratios
are the fragile estimands of any analysis, because a denominator with any mass near zero sends the
ratio's tails to infinity. Across the same \nbTwoNSeeds{} samples $g^\star$ carries a bias of
\nbTwoSeedGstarBias{} against a planted \nbTwoTrueGstar, a standard deviation of
\nbTwoSeedGstarSd, and its 94\,\% HDI covers the truth on \nbTwoSeedGstarCov{} of \nbTwoNSeeds.
That is the number to have in hand \emph{before} reading a cutoff off one sample, and
\cref{sec:seg:euro} reads one.

\subsection*{Falsifying the graph, and what that cannot do}

On real data, recovery-against-truth is unavailable and something must take its place. The standard
move is to test the conditional independences the DAG \emph{implies} against the data: a wrong graph
shows up as an independence the graph predicts and the data violate. Here \nbTwoDagTests{} implied
independences are tested at the 5\,\% level and \nbTwoDagViolations{} are violated --- the graph is
consistent with the data.

Hold that result lightly. \Cref{sec:seg:wrong} refits a model that is \emph{wrong} in a way that
destroys every individual-level deliverable in the chapter, and this same test passes it without
complaint.

\subsection*{A model-free second opinion}

\Cref{as:seg:form} is the assumption the simulator quietly granted. The cross-check that survives on
real data is a learner that commits to no functional form at all: a BART T-learner
\citep{chipman2010, kunzel2019}, which fits the two spend surfaces --- emailed and not --- with
sums of regression trees and differences them. Near mean engagement it returns \eur{\nbTwoBartLow}
for the low tier and \eur{\nbTwoBartHigh} for the high tier, against the structural model's
\eur{\nbTwoCateLow} and \eur{\nbTwoCateHigh}. A nonparametric learner lands on the same segment
effects: the structure is not doing the work, the data are.

% =====================================================================================
\section{Point estimate versus posterior}
\label{sec:seg:compare}

\input{build/tables/nb02_compare}

This is the section the chapter was written for, and it is the one where the machinery loses.

\subsection*{They agree, rung for rung}

\Cref{tab:seg:compare} puts every rung of the bracket on the same estimand, twice: once classically,
once Bayesianly. Read down the pairs. Classical no pooling and Bayesian no pooling: the same
numbers. Classical complete pooling and Bayesian complete pooling: the same numbers, wrong in the
same way, for both segments. Classical interaction OLS and the Bayesian interaction model this
chapter ships: the same numbers again. \textbf{The largest gap between the two members of any pair
is \eur{\nbTwoMaxPairGap}}, on an effect of \eur{\nbTwoTrueLow}--\eur{\nbTwoTrueHigh}.

This is not a disappointment; it is the lesson, and it has a name. With a well-specified likelihood
and enough data the posterior concentrates where the maximum-likelihood estimate does --- the
Bernstein--von Mises phenomenon \citep{gelman2013, kleijn2012} --- and the two paradigms return the
same number. \textbf{The causal work was done by the randomization of \cref{as:seg:unconf}, not by
the machinery.} Once identification is granted, a Bayesian regression with weak priors on
\nbTwoN{} rows is doing arithmetic that a pair of OLS fits had already done. Anyone hoping that
``going Bayesian'' would move the point estimate should read this table as the correction: with data
this plentiful and priors this weak, it must not, and it does not.

\subsection*{What Bayes did not buy}

The list is longer than the brochure admits, and each item is measured rather than conceded.

\begin{enumerate}[leftmargin=*, itemsep=.35em]
  \item \textbf{Not precision.} The low tier's 90\,\% width is \eur{\nbTwoWidthNpClsLow} from the
        crudest classical arm (no pooling), \eur{\nbTwoWidthIntClsLow} from the classical
        interaction OLS, and \eur{\nbTwoWidthIntBayesLow} from the Bayesian interaction model. The
        three land on top of each other --- a ratio of \nbTwoWidthRatio$\times$ against the crudest
        classical arm. At this $n$, sharing the nuisance parameters buys approximately nothing: the
        segments are fat and the effect is loud.
  \item \textbf{Not accuracy.} On this draw the low tier runs \eur{\nbTwoSampleErrLow} high in \emph{every} arm
        that estimates it separately --- classical and Bayesian, split and interacted, all four
        rows --- because the error is in the sample. (The two complete-pooling rows miss by
        \eur{\nbTwoCpMissLow}, but for an unrelated reason: they are not estimating this segment at
        all. They are reporting the ATE.)
  \item \textbf{Not the heterogeneity verdict.} The classical interaction coefficient --- the
        frequentist $\beta_3$, with an HC1 standard error --- is \eur{\nbTwoCgap} with a 90\,\%
        confidence interval of $[\eur{\nbTwoCgapLo}, \eur{\nbTwoCgapHi}]$, which excludes zero. A
        frequentist reaches ``the segments truly differ'' on the same data with no posterior at all.
        $\Prob(\beta_3 > 0) = \nbTwoPGap$ and ``the confidence interval excludes zero'' are two ways
        of reporting the same overwhelming evidence, in two vocabularies, and the honest analyst
        reports whichever the reader understands.
  \item \textbf{Not the pooling win.} \Cref{tab:seg:sweep} already settled this: at this chapter's
        segment sizes the learned $\lambda$ is \nbTwoRealLam, so borrowing strength is worth almost
        nothing. Partial pooling is a \emph{conditional} win, and this chapter is not the condition.
\end{enumerate}

Where the Bayesian model does pull ahead of the classical arms, it does so on \textbf{coherence},
not width: one fit produces the two segment CATEs, the posterior of their gap, the continuous
engagement profile, and the break-even engagement $g^\star = (c - \beta_1)/\beta_4$ --- a ratio of
estimated coefficients whose uncertainty falls straight out of the joint posterior, while the
classical route needs the delta method or \citeauthor{fieller1954}'s theorem \citep{fieller1954} and
a warning label. Coherence is a real benefit. It is a \emph{modest} one, and pretending otherwise
would cost this chapter its credibility.

\subsection*{What only the posterior can do}

One thing, and it is the thing the business asked for. \Cref{tab:seg:decide} reports
$\Prob(\tau_s > c)$ for each segment. The classical column of \cref{tab:seg:compare} reports that
the low tier's interval straddles the \eur{\nbTwoCost} cost, so the test \emph{cannot reject
break-even}.

Now try to run a campaign on that sentence. Do we mail them? ``Cannot reject'' is a statement about
a procedure's behaviour under a hypothesis. It is exactly as compatible with a segment that is
20\,\% likely to clear its cost as with one that is 60\,\% likely, and it will not say which,
although those two worlds call for opposite campaigns. The posterior says
$\Prob(\tau_{\mathrm{low}} > c) = \nbTwoPLowGtCost$ and
$\Prob(\tau_{\mathrm{high}} > c) = \nbTwoPHighGtCost$. Those are numbers, and
\cref{sec:seg:euro}'s send rule is written in exactly that currency.

\begin{remark}[The honest ledger]
Bayes did not rescue the estimate here --- it was never in danger. Two OLS fits reached the same two
numbers, with the same two interval widths, and the same verdict on the gap. What Bayes bought is
the right to make a probability statement \emph{about the effect}, which is the statement the send
rule is written in, and the ability to carry the full shape of the uncertainty --- rather than two
interval endpoints --- into the break-even engagement, the profit ladder and the cost sweep. That is
the whole win. On a borderline segment it is worth precisely the difference between mailing
\pc{\nbTwoShareLow} of the base and not.
\end{remark}

% =====================================================================================
\section{The decision, in euros}
\label{sec:seg:euro}

\subsection*{The send rule}

The rule is a stated convention, applied to the posterior of each segment's effect at a contact cost
of $c = \eur{\nbTwoCost}$:
%
\begin{equation}
  \Prob(\tau_s > c)\ \begin{cases}
    > 0.8 & \Rightarrow\ \textsc{send},\\
    < 0.5 & \Rightarrow\ \textsc{skip},\\
    \text{otherwise} & \Rightarrow\ \textsc{test further}.
  \end{cases}
  \label{eq:seg:rule}
\end{equation}
%
Note what \cref{eq:seg:rule} does \emph{not} say: it does not send wherever the mean beats the cost.
A segment whose posterior mean clears \eur{\nbTwoCost} by a whisker, with half its mass below the
line, is a coin flip dressed as a decision. \Cref{tab:seg:decide} applies the rule. High-value
customers are a \nbTwoActionHigh{} at $\Prob = \nbTwoPHighGtCost$; low-value customers are a
\nbTwoActionLow{} at $\Prob = \nbTwoPLowGtCost$ --- and it is precisely on that hesitant verdict that
the classical arm had nothing to offer but ``cannot reject''.

\input{build/tables/nb02_decide}

\subsection*{The break-even engagement}

A segment verdict is a blunt instrument when the effect within the segment slopes. The model's own
answer is sharper: solve $\tau(x) = c$ for engagement, using \cref{eq:seg:tauhat},
%
\begin{equation}
  g^{\star}_s \;=\; \frac{c \;-\; \beta_1 \;-\; \beta_3\, v_s}{\beta_4} ,
  \label{eq:seg:gstar}
\end{equation}
%
which turns a SKIP into a \emph{conditional} send. For the low tier the posterior of
\cref{eq:seg:gstar} gives $g^\star = \nbTwoGstarLow$ (94\,\% HDI
$[\nbTwoGstarLowLo, \nbTwoGstarLowHi]$), against a planted \nbTwoTrueGstar: mail a low-value customer
only once they are more engaged than that. For the high tier $g^\star = \nbTwoGstarHigh$ --- below
zero, i.e.\ off the scale --- so high-value customers clear the \eur{\nbTwoCost} cost at
\emph{every} engagement level. Always send.

\Cref{eq:seg:gstar} is where the coherence of the Bayesian fit finally pays a dividend the classical
arm cannot match, and also where it is at its most fragile, and both statements must be made
together. The posterior of a \emph{ratio} inherits amplified uncertainty near a crossing, and on this
sample the fitted cutoff sits below the true \nbTwoTrueGstar{} because the low tier's CATE runs
\eur{\nbTwoSampleErrLow} high and the ratio magnifies that error. \Cref{sec:seg:valid}'s multi-seed check is what keeps
this honest: across fresh samples $g^\star$ is centred (bias \nbTwoSeedGstarBias) and its interval
covers at \nbTwoSeedGstarCov{} of \nbTwoNSeeds. Read the printed band as an honest starting range for
a cutoff to be confirmed on a holdout --- not as a guarantee, and never as a point.

\subsection*{What targeting is worth}

The question that justifies the whole analysis is not statistical. It is: \emph{how much money does
segmenting make, against treating everyone as the average?} Because the per-customer truth is known
here --- and only a simulator can settle this --- four policies can be scored against it, forming a
ladder of increasing resolution: pooled (one all-or-nothing decision on the ATE), segment rules (one
decision per tier), individual (a decision per customer on $\hat\tau(x)$), and oracle (the unbeatable
rule that knows every true effect). The gaps between the rungs are, literally, the euros that finer
targeting is worth (\cref{tab:seg:policy}).

\input{build/tables/nb02_policy}

The headline --- \eur{\nbTwoValueTargeting} for individual targeting over pooling on this
\nbTwoN-customer sample, about \eur{\nbTwoValuePerHundredk} per 100{,}000 customers --- is itself an
estimate, and the same discipline demanded of every coefficient is demanded of it. Recomputing the
whole policy ladder \emph{per posterior draw} (each draw implies its own $\hat\tau(x)$, hence its own
send rules, hence its own realised profit) gives a 94\,\% HDI of
$[\eur{\nbTwoValueLo}, \eur{\nbTwoValueHi}]$ and $\Prob(\text{value} > 0) = \nbTwoValuePGtZero$. The
value of targeting is sign-certain, which is the claim a CMO needs, and it is bounded, which is the
claim a CFO needs.

\subsection*{Where the pooled policy snaps}

\Cref{fig:seg:costsweep} sweeps the contact cost and re-runs all four policies at each price. At a
low cost every policy earns the same, because ``email everyone'' is optimal and no analysis is
needed. As the cost rises the policies peel apart at their own decision thresholds --- and the
pooled rule fails in a specific, ugly way. It keeps mailing \emph{everybody} until the cost passes
the \emph{pooled} ATE of \eur{\nbTwoAte}, and then collapses to mailing \emph{nobody}, surrendering
in one step all the high-value profit that segment and individual targeting go on earning. The
segment rule, by contrast, drops the low tier once the cost passes \eur{\nbTwoCateLow} and keeps the
high tier.

\input{build/floats/nb02_costsweep}

That brittleness is the cost of one-size-fits-all, and it is not a rounding error in a spreadsheet;
it is a policy with no middle setting. The euro consequence of \cref{sec:seg:decision}'s bracket is
visible in a single chart.

% =====================================================================================
\section{What can go wrong}
\label{sec:seg:wrong}

\subsection*{Subgroup fishing --- and the cure, measured}

The most famous failure of segment analysis is not a bias in any estimator. It is
\textbf{multiplicity}. Slice a campaign twenty ways and something will look responsive by chance;
the analyst then writes a confident slide about ``segment 13''.

The notebook stages the trap on its own data (\cref{fig:seg:fishing}): the real outcome, the real
randomized email, and \nbTwoNFish{} coin-flip ``segments'' with zero true moderation by construction.
Each individual 95\,\% test is honest. The batch is not: \nbTwoNFish{} independent nulls at the
5\,\% level yield about \nbTwoFishExpected{} false discovery on average, and this run produces
\nbTwoFishHits. A responsive segment, manufactured from a coin flip.

\input{build/floats/nb02_fishing}

The defences, in order of cheapness. \emph{Pre-register} the segments to be acted on, before the
campaign runs --- which is what this chapter's own analysis did, with two moderators chosen for
business reasons, and that design choice protects more than any model can. \emph{Report posterior
probabilities} rather than one-shot significance stars. Control the false discovery rate if a scan
is genuinely unavoidable \citep{benjamini1995}. And --- the defence this chapter can now point at
rather than promise --- \emph{shrink}. Pointed at these very labels, \cref{eq:seg:shrink} returns
$\lambda = \nbTwoNoiseLam$ and collapses all \nbTwoNNoise{} sham segments onto the grand mean
(\cref{fig:seg:shrink}). Fishing dies when noisy estimates are not allowed to stand alone.

\subsection*{Interactions whisper; averages shout}

The mirror-image failure is that \emph{real} interactions are intrinsically hard to detect. In a
balanced test split into two equal segments, the gap is a difference of two segment effects, so its
sampling variance is the \emph{sum} of theirs --- and each is estimated from half the sample:
%
\begin{equation}
  \Var(\hat\beta_3) \;=\; \Var(\hat\tau_{\mathrm{high}}) + \Var(\hat\tau_{\mathrm{low}})
    \;\approx\; 2\,\Var(\hat\tau) + 2\,\Var(\hat\tau) \;=\; 4\,\Var(\hat\tau),
  \label{eq:seg:power}
\end{equation}
%
where $\hat\tau$ is the ATE estimate on the same $n$. The gap's standard error is therefore about
twice the ATE's --- verified on this chapter's own posterior, where the standard deviation of the
gap is \eur{\nbTwoGapSd} against \eur{\nbTwoAteSd} for the ATE, a ratio of
\nbTwoGapSdRatio$\times$ against a theoretical 2. Required sample grows with
$(\mathrm{se}/\text{effect})^2$, so an interaction \emph{half} the size of the main effect, carrying
twice its standard error, needs $(2 / \tfrac12)^2 = 16$ times the rows for equal power. This is
\citeauthor{gelman2018interactions}'s rule of thumb \citep{gelman2018interactions,
gelmanhill2007}, and its operational content is blunt: \textbf{an A/B test sized to measure the
average is roughly sixteen times too small to measure who it works on.}

The gap is detected easily here only because the planted gap is huge. Resolving the \emph{borderline}
call --- the low tier's effect against the \eur{\nbTwoCost} cost --- is the hard problem, and
\cref{eq:seg:rule}'s \textsc{test further} verdict must not end as a shrug. If the truth sits at
today's posterior mean, the sample that would place 95\,\% of the low-value posterior on one side of
the cost line is \nbTwoNFollowup{} customers. ``Run a bigger test'' now has a number attached, and
the closer the effect sits to the cost, the more brutally that number grows.

\subsection*{The functional form, which no graph test can check}

\Cref{as:seg:form} was handed to the model by the simulator: the true moderation was linear and the
specification said so. On real data nobody hands the analyst the form. So the notebook breaks the
model deliberately, refitting with the \code{email:engagement} moderation \emph{omitted} --- the
guess ``engagement predicts spend but does not moderate the email'', which would sound perfectly
reasonable before seeing the truth.

\input{build/floats/nb02_misspec}

Three things happen (\cref{fig:seg:misspec}), and the pattern is the lesson.

\begin{enumerate}[leftmargin=*, itemsep=.35em]
  \item \textbf{The segment-level answers survive}: \eur{\nbTwoBadLow} and \eur{\nbTwoBadHigh},
        against the correctly-specified \eur{\nbTwoCateLow} and \eur{\nbTwoCateHigh}. Omitting the
        moderation folds its average contribution into the baseline email coefficient, so at the
        segment level the error largely cancels. An analyst checking only the segment table would
        see nothing wrong.
  \item \textbf{Everything individual breaks.} Per-customer recovery degrades from RMSE
        \eur{\nbTwoRmseInd} to \eur{\nbTwoRmseBad}; the moderation profile is flat where the truth
        slopes, so the model now over-rates disengaged customers and under-rates engaged ones in
        \emph{both} tiers; and \cref{eq:seg:gstar} --- the chapter's operational deliverable ---
        cannot even be written down, because the wrong model has no $\beta_4$.
  \item \textbf{The DAG test stays silent}: \nbTwoDagBadViolations{} violations, exactly as on the
        correct model. This is not a defect of the test but a fact about what it tests. The wrong
        model has the \emph{same graph} as the right one --- the same nodes, the same arrows; only
        the algebra on the email$\rightarrow$spend arrow changed, and no conditional-independence
        test can see that. \textbf{Graph checks falsify your arrows. They cannot falsify your
        algebra.}
\end{enumerate}

The tools for the algebra are the ones \cref{sec:seg:valid} ran: the posterior predictive check,
multi-seed recovery where a simulator exists, and --- on real data, where it does not --- a
flexible model-free cross-check. Run all three, or accept that the functional form is an untested
assumption on which every individual-level number depends.

\subsection*{The limits}

Three remain, and they are properties of the problem rather than defects of the code.
\emph{Heterogeneity is only as good as the moderator that was measured}: if the real driver of
differential response is not the value tier or the engagement score, no interaction will catch it.
\emph{Interaction is not causation of the moderator}: nothing here claims that \emph{making} a
customer high-value would change their responsiveness, only that responsiveness differs across
pre-existing tiers. And \emph{the randomization is doing all the causal work}: observationally --- an
email sent to whoever the CRM liked --- a hidden common cause of email and spend would bias every
segment effect in the chapter, and no amount of pooling, shrinking or sampling would repair it.

% =====================================================================================
\section{Takeaways}
\label{sec:seg:take}

\begin{enumerate}[leftmargin=*, itemsep=.45em]
  \item \textbf{The average is a policy, and it is usually the wrong one.} Complete pooling returns
        the tightest interval in the chapter (\pc{\nbTwoCpNarrowerPct} narrower than no pooling) and
        excludes both planted truths, missing the low tier by \eur{\nbTwoCpMissLow}. It is an
        excellent estimate of a quantity nobody asked for. At the \eur{\nbTwoCost} contact cost it
        mails \pc{\nbTwoShareLow} of the base at a loss --- and the cost sweep
        (\cref{fig:seg:costsweep}) shows it later abandoning the profitable tier as well.
        \textbf{Precision is not accuracy}, and precision around the wrong number is the more
        dangerous failure, because it looks like confidence.

  \item \textbf{Partial pooling is the interpolation, and $\lambda$ is a signal-to-noise ratio.}
        \Cref{eq:seg:precision} is a precision-weighted average of a segment's own estimate and the
        grand mean --- not a fudge, and not Bayesian: it is \citeauthor{stein1956}'s
        \citeyearpar{stein1956}, and \citet{efron1975} showed it working on batting averages half a
        century ago. The two classical arms are its two limits.

  \item \textbf{The shrinkage is adaptive, and that is what makes it a discovery procedure.}
        Pointed at \nbTwoNNoise{} coin-flip segments, it recovers
        $\hat\sigma_\tau = \eur{\nbTwoNoiseSigma}$, sets $\lambda = \nbTwoNoiseLam$, and collapses
        every sham ``responsive segment'' onto the grand mean: \emph{partial pooling discovers that
        there is no heterogeneity to find}. No-pooling cannot ever reach that conclusion --- it will
        hand back as many numbers as there are labels, and some will look significant
        (\cref{fig:seg:fishing}). Pointed at the two real tiers, the same estimator returns
        $\lambda = \nbTwoRealLam$ and declines to shrink.

  \item \textbf{And it is worth little here --- say so.} At this chapter's segment sizes the learned
        $\lambda$ is near 1, and shrinkage beats plain no pooling by \eur{\nbTwoSweepThinSh} against
        \eur{\nbTwoSweepThinNp} even at the thinnest cells. Partial pooling is a \emph{conditional}
        win: its value scales with how noisy the segments are relative to how different they truly
        are. A book that sells it as a free precision win is selling something.

  \item \textbf{On this problem the Bayesian layer bought nothing on the number.} Classical and
        Bayesian agree rung for rung, to \eur{\nbTwoMaxPairGap} on a
        \eur{\nbTwoTrueLow}--\eur{\nbTwoTrueHigh} effect
        (\cref{tab:seg:compare}) --- Bernstein--von Mises, on schedule. It bought nothing on interval
        width (\nbTwoWidthRatio$\times$ the crudest classical arm) and nothing on the heterogeneity
        verdict (the classical $\beta_3$ interval, $[\eur{\nbTwoCgapLo}, \eur{\nbTwoCgapHi}]$,
        excludes zero on its own). What it bought is \emph{coherence} --- one fit yielding both
        effects, their gap, the continuous profile and the ratio estimand $g^\star$ --- and the
        decision quantity.

  \item \textbf{The decision quantity is the whole win.} $\Prob(\tau_s > c)$ has no classical
        counterpart. ``Cannot reject break-even'' is as compatible with a 20\,\% chance of clearing
        cost as with a 60\,\% chance, and those two worlds call for opposite campaigns.
        \Cref{eq:seg:rule} consumes nothing else, and neither does the CMO.

  \item \textbf{Robust standard errors are insurance, not magic.} HC1 and the iid formula differ by
        \pc{\nbTwoSeGapPct} here, because the simulator draws homoskedastic noise. The premium is a
        rounding error --- and on real spend data, where high-value customers vary as well as spend
        more, it would not be. Pay a premium that is zero when you do not need it.

  \item \textbf{Graph checks falsify your arrows, never your algebra.} A model with the moderation
        omitted passes the DAG-falsification test (\nbTwoDagBadViolations{} violations), keeps the
        segment answers nearly intact, and destroys every individual-level deliverable in the
        chapter --- the per-customer RMSE, the moderation profile, and the break-even rule, which it
        cannot even express. The functional form is an assumption, and it needs a PPC, a multi-seed
        recovery check and a model-free cross-check, not a graph.

  \item \textbf{Budget for the interaction, not for the average.} The gap's standard error is
        \nbTwoGapSdRatio$\times$ the ATE's on this very posterior, and an interaction half the size
        of a main effect needs sixteen times the sample for equal power
        \citep{gelman2018interactions}. Interactions whisper; averages shout. An experiment sized to
        measure the average is not sized to measure who it works on --- and the follow-up that would
        settle this chapter's borderline tier needs \nbTwoNFollowup{} customers.
\end{enumerate}

\notebooknote{%
  This chapter is \code{notebooks/02\_segment\_effects.ipynb}. It runs in a fast teaching mode
  (\code{CMP\_FAST=1}, short chains) and a full mode (\code{CMP\_FAST=0}); every number in this
  chapter comes from the full run, emitted by \code{cmp.report} and injected here as a macro. The
  simulator is \code{cmp.dgp.segment\_customers}, the classical arms are \code{cmp.classical}
  (\code{no\_pooling}, \code{complete\_pooling}, \code{ols}), the structural model is fitted with
  \code{pathmc}, and the seed is fixed at \nbTwoSeed. \Citet{gelmanhill2007} is the standard
  practitioner's treatment of the pooling bracket; \citet{efron1977} is the shortest and best
  introduction to why shrinkage works at all.}
